% ============================================================
% PHẦN 3.12 — CHỈ SỐ ĐO LƯỜNG HIỆU QUẢ
% ------------------------------------------------------------
% Thuộc: PA3-05
% Người phụ trách chính: Nguyễn Công Toàn – Finance Lead
% Phối hợp dữ liệu: Lê Nguyên Nhật
%
% CÂU HỎI CẦN TRẢ LỜI:
%   - KPI nào đo nhận biết?
%   - KPI nào đo mức quan tâm?
%   - KPI nào đo chuyển đổi?
%   - KPI nào đo kích hoạt và sử dụng?
%   - KPI nào đo mức độ hài lòng?
%   - KPI nào quyết định tiếp tục hoặc dừng một kênh?
%   - Dữ liệu được lấy từ đâu?
%   - Ai cập nhật số liệu?
%   - Tần suất cập nhật là hàng ngày, hàng tuần hay sau mỗi pilot?
%
% OUTPUT BẮT BUỘC:
%   1. Bảng KPI đầy đủ:
%      KPI | Ý nghĩa | Công thức | Nguồn dữ liệu | Tần suất | Người phụ trách | Mục tiêu
%
%   KPI BẮT BUỘC theo đề bài:
%      - Số người tiếp cận
%      - Lượt tương tác
%      - Lượt đăng ký dùng thử
%      - Tỷ lệ chuyển đổi từ người quan tâm sang người dùng
%      - Chi phí có một khách hàng mới (CAC)
%      - Phản hồi của khách hàng
%
%   KPI B2B NÊN BỔ SUNG:
%      - Số khách hàng phù hợp ICP
%      - Số phản hồi từ tiếp cận trực tiếp
%      - Số lịch demo
%      - Tỷ lệ tham gia demo
%      - Số đơn vị đồng ý pilot
%      - Tỷ lệ hoàn thành onboarding
%      - Thời gian từ đăng ký đến phiên đầu tiên
%      - Số lead được ghi nhận trong pilot
%      - Tỷ lệ pilot chuyển sang trả phí
%      - Mức độ hài lòng
%      - Ý định tiếp tục sử dụng hoặc giới thiệu
%
%   2. Công thức phễu cơ bản:
%      - Tỷ lệ phản hồi = Số khách hàng phản hồi / Số khách hàng tiếp cận × 100%
%      - Tỷ lệ đặt lịch demo = Số lịch demo / Số khách hàng phản hồi × 100%
%      - Tỷ lệ chuyển đổi pilot = Số khách hàng tham gia pilot / Số đã demo × 100%
%      - CAC = Tổng chi phí bán hàng và quảng bá / Số khách hàng trả phí mới
%
%   3. Bảng điều kiện tiếp tục hoặc dừng kênh:
%      Tiếp tục / Điều chỉnh / Tạm dừng / Loại bỏ
%
% TIÊU CHÍ HOÀN THÀNH:
%   - Mỗi KPI phải có công thức và nguồn dữ liệu
%   - Không đặt mục tiêu phần trăm tùy ý mà không giải thích
%   - Phân biệt: Lead / Lead phù hợp ICP / Người đăng ký demo /
%     Người dùng thử / Khách hàng trả phí
%   - Khi chưa có dữ liệu thật → ghi là "mục tiêu thử nghiệm" hoặc "giả định ban đầu"
% ============================================================

\subsection{Chỉ số đo lường hiệu quả}

% TODO (Nguyễn Công Toàn + Lê Nguyên Nhật): Viết nội dung mục 3.12
% Toàn phụ trách bảng KPI và công thức
% Nhật phụ trách nguồn dữ liệu kỹ thuật và danh sách sự kiện tracking (liên kết 3.8)
